%%% FIELD proof-large-k-honest-length
Write
\[
\mathbb Q_{\mathbf P}^{N}
:=\bigotimes_{k=1}^{K}P_k^{\otimes m}
\quad\text{on}\quad
(\mathbb B_d)^{Km}.
\]
This is the observable law supplied by Assumption~\texttt{ass:product-experiment}; Assumption~\texttt{ass:balanced-array} gives \(Km=N\).  The inferential target in this theorem is the second-moment endpoint \(L(\mathbf P)\), not the centered quantity \(V(\mathbf P)\).  Theorem~\texttt{thm:sharp-residual-envelope} and the positivity of \(K^2\) give the deterministic sandwich
\[
0\le K^2\{L(\mathbf P)-\|\bar\mu\|_2^2\}\le r,
\qquad
\|\bar\mu\|_2^2\le L(\mathbf P)
\le \|\bar\mu\|_2^2+\frac r{K^2}.                                      \tag{1}
\]

We first prove uniform honesty by a Hilbert-space second-moment calculation followed by a pointwise Markov inequality.  Put
\[
Y_{ki}:=X_{ki}-\mu_k,
\qquad
\varepsilon_N:=(\alpha N)^{-1/2}.
\]
The radius is well defined because \(N>0\) and \(\alpha\in(0,1/4)\).  Under Assumption~\texttt{ass:product-experiment}, the \(Y_{ki}\) are independent and have mean zero.  Since \(X_{ki}\in\mathbb B_d\), the variance identity and the support bound imply
\[
\mathbb E_{\mathbf P}\|Y_{ki}\|_2^2
=\mathbb E_{\mathbf P}\|X_{ki}\|_2^2-\|\mu_k\|_2^2
\le 1.                                                                    \tag{2}
\]
For two distinct sample indices, independence and centering give
\(
\mathbb E_{\mathbf P}\langle Y_{ki},Y_{\ell j}\rangle=0
\).
Consequently, using \(N=Km\),
\[
\begin{aligned}
\mathbb E_{\mathbf P}\|\bar X-\bar\mu\|_2^2
&=\frac1{N^2}
  \mathbb E_{\mathbf P}
  \left\|\sum_{k=1}^{K}\sum_{i=1}^{m}Y_{ki}\right\|_2^2 \\
&=\frac1{N^2}\sum_{k=1}^{K}\sum_{i=1}^{m}
  \mathbb E_{\mathbf P}\|Y_{ki}\|_2^2
\le \frac1N.
\end{aligned}                                                             \tag{3}
\]
The pointwise inequality
\[
\mathbf 1\{\|\bar X-\bar\mu\|_2>\varepsilon_N\}
\le \frac{\|\bar X-\bar\mu\|_2^2}{\varepsilon_N^2}
\]
and (3) yield, uniformly in \(\mathbf P\),
\[
\mathbb Q_{\mathbf P}^{N}
\{\|\bar X-\bar\mu\|_2\le\varepsilon_N\}
\ge 1-\frac{N^{-1}}{\varepsilon_N^2}=1-\alpha.                            \tag{4}
\]
On the event in (4), the reverse triangle inequality gives
\[
\max\{0,\|\bar X\|_2-\varepsilon_N\}
\le \|\bar\mu\|_2.
\]
After squaring both nonnegative sides and applying the lower inequality in (1), Definition~\texttt{def:mean-lower-bound} gives
\[
B_N^{\mathrm{mean}}
=\bigl(\max\{0,\|\bar X\|_2-\varepsilon_N\}\bigr)^2
\le \|\bar\mu\|_2^2
\le L(\mathbf P).                                                         \tag{5}
\]
Equations (4)--(5), followed by taking the infimum over \(\mathbf P\), prove
\[
\inf_{\mathbf P\in\mathcal P(\mathbb B_d)^K}
\mathbb Q_{\mathbf P}^{N}
\{B_N^{\mathrm{mean}}\le L(\mathbf P)\}
\ge1-\alpha.
\]
Thus \(B_N^{\mathrm{mean}}\in\mathcal H_\alpha\) by Definition~\texttt{def:honest-bounds}.

We next bound its expected excess length by deterministic interpolation between the endpoint sandwich (1) and the sample-mean error.  Let
\[
u:=\|\bar\mu\|_2,
\qquad
t:=\|\bar X\|_2,
\qquad
v:=(t-\varepsilon_N)_+.
\]
Convexity of \(\mathbb B_d\) implies \(u,t\in[0,1]\), and hence \(v\in[0,1]\).  The reverse triangle inequality and \(0\le t-v\le\varepsilon_N\) show
\[
|u-v|
\le |u-t|+|t-v|
\le \|\bar X-\bar\mu\|_2+\varepsilon_N.                                  \tag{6}
\]
For \(u,v\in[0,1]\), factorization gives
\[
|u^2-v^2|=|u-v|(u+v)\le2|u-v|.                                            \tag{7}
\]
Because \(L(\mathbf P)-u^2\ge0\), the elementary inequality
\((a+b)_+\le a+|b|\) for \(a\ge0\), together with (1), (6), and (7), yields
\[
\begin{aligned}
\{L(\mathbf P)-B_N^{\mathrm{mean}}\}_+
&=\{(L(\mathbf P)-u^2)+(u^2-v^2)\}_+ \\
&\le L(\mathbf P)-u^2+|u^2-v^2| \\
&\le \frac r{K^2}
   +2\|\bar X-\bar\mu\|_2+2\varepsilon_N.
\end{aligned}                                                             \tag{8}
\]
Cauchy--Schwarz applied to (3) gives
\(
\mathbb E_{\mathbf P}\|\bar X-\bar\mu\|_2\le N^{-1/2}
\).
Taking expectations in (8) therefore proves, uniformly in \(\mathbf P\),
\[
\mathbb E_{\mathbf P}
  [\{L(\mathbf P)-B_N^{\mathrm{mean}}\}_+]
\le \frac r{K^2}
 +(2+2\alpha^{-1/2})N^{-1/2}.                                             \tag{9}
\]
Since the displayed statistic is an admissible member of \(\mathcal H_\alpha\), Definition~\texttt{def:minimax-excess-length} and (9) imply
\[
E_{d,\alpha}(N,K)
\le C_{d,\alpha}\left(N^{-1/2}+\frac r{K^2}\right),
\qquad
C_{d,\alpha}:=\max\{1,2+2\alpha^{-1/2}\}.                               \tag{10}
\]

It remains to prove the all-procedure converse.  Use the ordered pair supplied by Lemma~\texttt{lem:bernoulli-two-point}, and abbreviate
\[
L_j:=L(\mathbf P^{(j)}),
\qquad
\mathbb Q_j:=\bigotimes_{k=1}^{K}(P_k^{(j)})^{\otimes m},
\qquad
\Delta:=L_1-L_0.
\]
That lemma gives
\[
\Delta\ge aN^{-1/2},
\qquad
\operatorname{TV}(\mathbb Q_0,\mathbb Q_1)\le\frac14.                   \tag{11}
\]
Let \(B_N\in\mathcal H_\alpha\) be arbitrary.  Uniform honesty at \(\mathbf P^{(0)}\) gives
\[
\mathbb Q_0\{B_N\le L_0\}\ge1-\alpha.                                  \tag{12}
\]
By the defining eventwise inequality for total variation, (11)--(12) imply
\[
\mathbb Q_1\{B_N\le L_0\}
\ge\mathbb Q_0\{B_N\le L_0\}
   -\operatorname{TV}(\mathbb Q_0,\mathbb Q_1)
\ge\frac34-\alpha
\ge\frac12,                                                              \tag{13}
\]
where the last inequality uses \(\alpha<1/4\).  On the event in (13), ordered separation gives
\(
(L_1-B_N)_+\ge L_1-L_0=\Delta
\).
Combining this fact with (11)--(13),
\[
\sup_{\mathbf P}
\mathbb E_{\mathbf P}[(L(\mathbf P)-B_N)_+]
\ge\mathbb E_{\mathbb Q_1}[(L_1-B_N)_+]
\ge\frac a2N^{-1/2}.                                                      \tag{14}
\]
Because \(B_N\) was arbitrary, taking the infimum over \(\mathcal H_\alpha\) in (14) proves
\[
E_{d,\alpha}(N,K)\ge c_{d,\alpha}N^{-1/2},
\qquad
c_{d,\alpha}:=\min\{a/2,1\}>0.                                          \tag{15}
\]
The choices in (10) and (15) satisfy
\(0<c_{d,\alpha}\le C_{d,\alpha}<\infty\),
and the constants depend at most on \(d\) and \(\alpha\).

Finally, if \(K\ge c_0N^{1/4}\), then \(r=\min\{K,d\}\le d\) and positivity of \(c_0\) gives
\[
\frac r{K^2}
\le \frac d{c_0^2N^{1/2}}.                                               \tag{16}
\]
Equations (10), (15), and (16) prove
\(E_{d,\alpha}(N,K)\asymp N^{-1/2}\)
on this regime.  Assumption~\texttt{ass:legal-triangular-array} ensures that this asymptotic assertion is read only along the declared balanced arrays with \(m=N/K\to\infty\); all inequalities above are finite-sample inequalities whenever the frozen conditions \(N=Km\), \(K\ge3\), and \(m\ge2\) hold.

The construction is independently operational: its formula uses only \(\bar X\), \(N\), and \(\alpha\), with the fully specified radius \((\alpha N)^{-1/2}\).  It neither evaluates the empirical endpoint \(\widehat L\), solves a multi-marginal optimal-transport problem, nor selects an existential empirical-process constant \(C_d\).  Thus the proof establishes the stated mean-only, no-transport computational baseline without identifying \(L(\mathbf P)\) with the centered residual \(V(\mathbf P)\).
